\section*{Capítulo \ref{ch:g5:water-shapes-the-land} --- El agua modela el terreno}

\begin{solution}{exo:g5:water-shapes-the-land:1}
(a) \omterm{def:g5:water-shapes-the-land:weathering}{Meteorización}: la \omterm{def:g3:stones-pebbles-sand:rock}{roca} se rompe donde está. (b) \omterm{def:g5:water-shapes-the-land:erosion}{Erosión}: la \omterm{def:g3:stones-pebbles-sand:sand}{arena} se
desgasta y se la lleva el agua. (c) \omterm{def:g5:water-shapes-the-land:deposition}{Sedimentación}: el \omterm{def:g3:stones-pebbles-sand:sand}{barro} se deposita.
(d) \omterm{def:g5:water-shapes-the-land:weathering}{Meteorización}: la raíz rompe la piedra donde está.
\end{solution}

\begin{solution}{exo:g5:water-shapes-the-land:2}
El \omterm{def:g5:water-shapes-the-land:source-mouth}{nacimiento} es el lugar donde empieza el \omterm{def:g1:landscapes-around-us:water}{río} (un manantial, un glaciar
que se funde, una tierra alta y húmeda). La \omterm{def:g5:water-shapes-the-land:source-mouth}{desembocadura} es el lugar donde
termina, en el \omterm{def:g2:water-on-earth:ocean}{mar}, en un \omterm{def:g1:landscapes-around-us:water}{lago} o en un \omterm{def:g1:landscapes-around-us:water}{río} más grande. Un \omterm{def:g5:water-shapes-the-land:source-mouth}{afluente} es un \omterm{def:g1:landscapes-around-us:water}{río}
más pequeño que desemboca en otro más grande.
\end{solution}

\begin{solution}{exo:g5:water-shapes-the-land:3}
El \omterm{def:g5:water-shapes-the-land:source-mouth}{nacimiento}; el \omterm{def:g1:landscapes-around-us:landform}{valle} estrecho en V; los \omterm{def:g5:water-shapes-the-land:meander}{meandros}; el \omterm{def:g5:water-shapes-the-land:delta}{delta}.
\end{solution}

\begin{solution}{exo:g5:water-shapes-the-land:4}
(a) Falso: la helada, el agua, las raíces, el calor y el frío rompen poco
a poco incluso la \omterm{def:g3:stones-pebbles-sand:rock}{roca} dura (\omterm{def:g5:water-shapes-the-land:weathering}{meteorización}). (b) Verdadero: cuanto más
rápida es el agua, más grandes son los trozos que puede mover. (c) Falso:
los trozos más pesados se depositan primero; el \omterm{def:g3:stones-pebbles-sand:sand}{barro} más fino se deposita
el último, donde el agua está casi quieta, en un \omterm{def:g1:landscapes-around-us:water}{lago} o en el \omterm{def:g2:water-on-earth:ocean}{mar}.
\end{solution}

\begin{solution}{exo:g5:water-shapes-the-land:5}
A está en la parte de fuera de la curva: allí el agua es rápida y profunda,
y desgasta la orilla (\omterm{def:g5:water-shapes-the-land:erosion}{erosión}). B está en la parte de dentro de la curva:
allí el agua es lenta y deja \omterm{def:g3:stones-pebbles-sand:sand}{arena} (\omterm{def:g5:water-shapes-the-land:deposition}{sedimentación}).
\end{solution}

\begin{solution}{exo:g5:water-shapes-the-land:6}
En las \omterm{def:g1:landscapes-around-us:landform}{montañas} el \omterm{def:g1:landscapes-around-us:water}{río} es empinado y rápido: puede hacer rodar piedras
grandes y se lleva los granos pequeños. Cerca del \omterm{def:g2:water-on-earth:ocean}{mar} el \omterm{def:g1:landscapes-around-us:water}{río} es lento: ya no
puede mover piedras grandes, solo \omterm{def:g3:stones-pebbles-sand:sand}{arena} y \omterm{def:g3:stones-pebbles-sand:sand}{barro}, y los deja allí. Además,
por el camino, las piedras se han roto y se han desgastado hasta formar
granos más pequeños.
\end{solution}

\begin{solution}{exo:g5:water-shapes-the-land:7}
(a) Curso alto: $600 \div 6 = 100$ kilómetros. Curso medio:
$600 \div 2 = 300$ kilómetros. (b) Curso bajo: $600 - 100 - 300 = 200$
kilómetros. Como $3 \times 200 = 600$, es $\frac{1}{3}$ del \omterm{def:g1:landscapes-around-us:water}{río}.
(Comprobación: $\frac{1}{6} + \frac{1}{2} + \frac{1}{3} = \frac{1}{6} +
\frac{3}{6} + \frac{2}{6} = \frac{6}{6}$.)
\end{solution}

\begin{solution}{exo:g5:water-shapes-the-land:8}
Las manchas claras son \omterm{def:g3:stones-pebbles-sand:sand}{arena} y \omterm{def:g3:stones-pebbles-sand:pebble}{grava} depositadas en la parte de dentro de
las curvas, donde el agua es lenta (\omterm{def:g5:water-shapes-the-land:deposition}{sedimentación}). Los \omterm{def:g1:landscapes-around-us:water}{lagos} curvos
oscuros son antiguas curvas del \omterm{def:g1:landscapes-around-us:water}{río}: dos curvas se acercaron tanto que el
\omterm{def:g1:landscapes-around-us:water}{río} rompió el cuello de tierra que las separaba, tomó el atajo y dejó atrás
la curva vieja convertida en \omterm{def:g1:landscapes-around-us:water}{lago}.
\end{solution}

\begin{solution}{exo:g5:water-shapes-the-land:9}
Los prados llanos que hay junto al \omterm{def:g1:landscapes-around-us:water}{río} son su \omterm{def:g5:water-shapes-the-land:flood}{llanura de inundación}: la
tierra que el \omterm{def:g1:landscapes-around-us:water}{río} cubre cuando se desborda. Son llanos precisamente porque
el \omterm{def:g1:landscapes-around-us:water}{río} los construyó con su propio \omterm{def:g3:stones-pebbles-sand:sand}{barro}. Algún día el \omterm{def:g1:landscapes-around-us:water}{río} volverá a
desbordarse y llenará de agua el colegio, y el agua rápida de una \omterm{def:g5:water-shapes-the-land:flood}{crecida}
es peligrosa. El colegio debería construirse en terreno alto, y los prados
deberían quedarse como campos, prados o un parque.
\end{solution}

\begin{solution}{exo:g5:water-shapes-the-land:10}
(a) Un año tiene 12 meses: $12$ millardos $\div\, 12 = 1$ millardo (mil
millones) de toneladas al mes, de media. (b) Mil millones son
$1\,000\,000\,000$; divididos entre 10, son $100\,000\,000$: cien millones
de viajes de camión.
\end{solution}

\begin{solution}{exo:g5:water-shapes-the-land:11}
Una historia posible. El grano forma parte de una \omterm{def:g3:stones-pebbles-sand:rock}{roca} en lo alto de una
\omterm{def:g1:landscapes-around-us:landform}{montaña}. La helada abre una grieta a su alrededor, noche tras noche, hasta
que se desprende un trozo de \omterm{def:g3:stones-pebbles-sand:rock}{roca}: eso es la \omterm{def:g5:water-shapes-the-land:weathering}{meteorización}. El trozo cae al
arroyo de abajo. En el curso alto el agua es rápida y la pendiente fuerte:
el agua hace rodar el trozo (\omterm{def:g5:water-shapes-the-land:erosion}{erosión}) y, al chocar con otras piedras, este
se rompe en \omterm{def:g3:stones-pebbles-sand:pebble}{cantos} y después en granos de \omterm{def:g3:stones-pebbles-sand:sand}{arena}. En el curso medio el \omterm{def:g1:landscapes-around-us:water}{río}
se frena; el grano descansa un tiempo en la playa de \omterm{def:g3:stones-pebbles-sand:sand}{arena} de la parte de
dentro de un \omterm{def:g5:water-shapes-the-land:meander}{meandro}, hasta que una \omterm{def:g5:water-shapes-the-land:flood}{crecida} vuelve a llevárselo. En el
curso bajo el \omterm{def:g1:landscapes-around-us:water}{río} es lento y ancho; el grano viaja con el agua hasta la
\omterm{def:g5:water-shapes-the-land:source-mouth}{desembocadura}. Allí el agua casi se detiene, y el grano se deposita con el
\omterm{def:g3:stones-pebbles-sand:sand}{barro} y la \omterm{def:g3:stones-pebbles-sand:sand}{arena} del \omterm{def:g5:water-shapes-the-land:delta}{delta}: eso es la \omterm{def:g5:water-shapes-the-land:deposition}{sedimentación}.
\end{solution}
